\documentclass[11pt]{article}
\usepackage[margin=1in]{geometry}
\usepackage{amsmath,amssymb,amsthm,microtype}
\usepackage[T1]{fontenc}
\usepackage{lmodern}
\usepackage[colorlinks=true,linkcolor=blue,citecolor=blue,urlcolor=blue]{hyperref}
\newtheorem{theorem}{Theorem}
\newtheorem{lemma}[theorem]{Lemma}
\newcommand{\R}{\mathbb R}
\newcommand{\Q}{\mathbb Q}
\newcommand{\Z}{\mathbb Z}
\newcommand{\F}{\mathcal F}
\DeclareMathOperator{\sig}{sign}
\title{Deep derived pure braids need not be 1-solvable}
\author{Alec Kriebel\\Independent researcher\\
\small\href{https://orcid.org/0009-0001-9320-500X}{ORCID: 0009-0001-9320-500X}}
\date{October 15, 2026\quad (version 1.0)}
\begin{document}
\maketitle
\begin{abstract}
For every $m\geq3$ and every finite $k\geq1$, the $k$th ordinary
derived subgroup of the pure braid group contains a braid whose standard
closure is not integrally smoothly $1$-solvable. Consequently
$P_m^{(n+2)}\not\subset\F_n^m$ for every integer $n\geq1$,
answering negatively the derived-series containment question in the
Harvey--Park--Ray Oberwolfach contribution.
The proof combines established signature-quasimorphism results with
amenable-quotient rigidity and a pointwise signature obstruction to
$1$-solvability. The general quasimorphism ingredients are prior results.
An accompanying exact verification package also preserves an explicit
104-crossing example in $P_3^{(3)}$.
\end{abstract}

\section{Statement and conventions}
Write $G^{(0)}=G$ and $G^{(j+1)}=[G^{(j)},G^{(j)}]$, with
$[x,y]=xyx^{-1}y^{-1}$. Let $P_m\subset B_m$ be the classical disk pure
braid group. Interpret a pure braid as an ordered smooth string link and
write $\F_n^m$ for the integral smooth solvable filtration defined by
standard closures, as in Otto \cite[Definition 2.3]{Otto}.
Containment here concerns the string-link concordance classes represented
by these braids. Harvey's contribution, joint with Park and Ray,
asks whether $P_m^{(n+2)}\subset\F_n^m$ for $n\geq1$
\cite[Question 3, p.3021]{HPR}.
The series is the ordinary derived series, rather than a power subgroup.

\begin{theorem}\label{thm:main}
For every $m\geq3$ and every finite integer $k\geq1$,
\[
 P_m^{(k)}\not\subset\F_1^m.
\]
In particular, $P_m^{(n+2)}\not\subset\F_n^m$ for every $m\geq3$
and integer $n\geq1$.
\end{theorem}

The witnesses may depend on $k$. No single braid in the intersection of
all derived subgroups is asserted. For $m=1,2$ the positive derived
subgroups are trivial, since $P_1$ is trivial and $P_2\cong\Z$.
The proof below is existential at arbitrary depth. The separate explicit
example in Section~\ref{sec:explicit} addresses $k=3$ only.

\section{Two signature inputs}
Let $s(\beta)=\sigma(\widehat\beta)$ be the ordinary signature of the
standard oriented braid closure, with positive trefoil signature $-2$.
This is a link signature, without a knot restriction.
Gambaudo--Ghys \cite[Theorem A and Remark 1]{GG} show that $s$ has bounded
defect on $B_3$. Its homogenization
\[
 q(\beta)=\lim_{r\to\infty}\frac{s(\beta^r)}r
\]
is a homogeneous quasimorphism; their discussion on p.575 also gives
$|q-s|\leq6$ on $B_3$. A homogeneous quasimorphism satisfies
$q(g^r)=rq(g)$ for integers $r$, and has finite defect
$D(q)=\sup_{g,h}|q(gh)-q(g)-q(h)|$.

\begin{lemma}\label{lem:pure}
The restriction $q|_{P_3}$ is not a homomorphism.
\end{lemma}
\begin{proof}
Put $x=\sigma_1$, $y=\sigma_2$, and
\[
 a=x^2,\qquad b=y^2,\qquad h=yx^2y^{-1},\qquad
 z=(xy)^3=\Delta_3^2.
\]
These four braids are pure. The Artin relation $xyx=yxy$ gives
\[
 ahb=x^2yx^2y=x(xyx)xy=x(yxy)xy=(xy)^3=z.
\]
By \cite[Remark 3, p.545]{GG}, the two-strand closure of $x^\ell$
has signature $1-\ell$ for $\ell\geq1$. Adding a split trivial strand
preserves signature. Conjugate braid closures are isotopic, and $x,y$
are conjugate, so
$q(a)=q(b)=q(h)=-2$.
The three-strand specialization of \cite[Lemma 4.1, pp.568--569]{GG}
is $s(\alpha\Delta_3^4)=s(\alpha)-8$.
Thus $s(z^{2r})=-8r$ for all positive $r$, whence $q(z)=-4$.
A homomorphism on $P_3$ would instead give $q(z)=-6$.
\end{proof}

The next input is stated with its hypotheses because we use the
\emph{pointwise} signature at $-1$, including at Alexander roots.
It is a standard signature obstruction; we supply the local-coefficient
argument and cite the precise surgery formula.

\begin{lemma}\label{lem:topology}
If an algebraically split oriented link $L$ admits an integral smooth
$1$-solution in the sense of \cite[Definition 2.3]{Otto}, then
$\sigma_L(-1)=0$.
\end{lemma}
\begin{proof}
Let $M_L=\partial W$ be zero surgery and let
$\{L_i,D_i\}_{i=1}^r$ be the framed embedded hyperbolic basis of
$H_2(W;\Z)$ supplied by the solution. Both sets of surface groups lie
in $\pi_1(W)^{(1)}$, and inclusion induces
$H_1(M_L;\Z)\cong H_1(W;\Z)\cong\Z^m$.
The character taking every meridian to $-1$ therefore extends over $W$;
denote its rank-one rational local system by $\Q_\alpha$.

In a finite cellular structure, each twisted boundary matrix has integral
entries obtained by evaluating group elements at $\pm1$. Modulo $2$ it
is the ordinary cellular boundary matrix. A minor nonzero over
$\mathbb F_2$ has odd determinant, so
$\operatorname{rank}_{\Q}d_j^\alpha\geq
\operatorname{rank}_{\mathbb F_2}d_j$ for each $j$. The chain dimensions
are equal, and hence
\[
 \dim_{\Q}H_2(W;\Q_\alpha)
 \leq\dim_{\mathbb F_2}H_2(W;\mathbb F_2)=2r.
\]
The final equality follows from integral universal coefficients:
$H_2(W;\Z)$ is free of rank $2r$ and $H_1(W;\Z)$ has no $2$-torsion.

The local system restricts trivially to each basis surface because its
character kills the commutator subgroup. Each surface therefore has a
local-coefficient fundamental class after a choice of trivialization.
Their twisted pairings are zero except for the prescribed dual pairs,
whose single intersections have coefficients $+1$ or $-1$.
Framings make the self-pairings zero. Rescaling one class in each pair
gives $r$ hyperbolic blocks. These $2r$ classes are independent; the
dimension bound makes them a full basis. Thus the twisted form, including
any possible ambient radical, has signature zero. The ordinary form is
also hyperbolic and has signature zero. Extension to real or complex
coefficients preserves these signatures.

Toffoli's bounding formula \cite[(1.1)]{Toffoli} gives
$\rho_\alpha(M_L)=\sig(W)-\sig_\alpha(W)=0$.
His single-colour surgery formula \cite[Remark 4.25, (4.21)]{Toffoli}
gives $\rho_\alpha(M_L)=-\sigma_L(-1)$: the entire framed linking
matrix is zero for this algebraically split, zero-framed surgery, so
all correction terms vanish. This pointwise formula applies at every
unit value other than $1$; no Alexander-root exclusion or average of
signatures is used. Therefore $\sigma_L(-1)=0$.
\end{proof}

\section{Persistence through the derived series}
We recall a consequence of established quasimorphism theory, giving a
proof to make the application self-contained apart from amenable rigidity.
Kawasaki--Kimura--Maruyama--Matsushita--Mimura
\cite[Lemma 7.12, p.1207]{KKMMM} prove the stronger identity
$D_G(f)=D_N(f|_N)$ whenever $G/N$ is solvable.
In particular, the following group-theoretic fact is not claimed new.

\begin{lemma}\label{lem:derived}
If a homogeneous real quasimorphism $f$ on $G$ is not a homomorphism,
then $f|_{G^{(k)}}$ is nonzero for every finite $k\geq1$.
\end{lemma}
\begin{proof}
Suppose $f$ vanishes on $N=G^{(k)}$. This subgroup is normal.
For $g\in G$, $u\in N$ and positive $r$, normality gives
$(gu)^r=g^r u_r$ with $u_r\in N$. Applying the defect bound once yields
\[
 |r f(gu)-r f(g)|\leq D(f).
\]
After division by $r$ and passage to infinity, $f(gu)=f(g)$.
Thus $f$ descends to a homogeneous quasimorphism on $G/N$.
This quotient is solvable, hence amenable. A homogeneous real
quasimorphism on an amenable group is a homomorphism
\cite[Proposition 2.65, p.33; p.27]{Calegari}.
Pulling it back contradicts the hypothesis.
\end{proof}

\begin{proof}[Proof of Theorem~\ref{thm:main}]
Fix finite $k\geq1$. Lemmas~\ref{lem:pure} and~\ref{lem:derived}
give $v\in P_3^{(k)}$ with $q(v)\ne0$.
Choose a positive $r$ with $r|q(v)|>6$. Since
$q(v^r)=rq(v)$ and $|q-s|\leq6$, the braid $\gamma=v^r$ satisfies
$s(\gamma)\ne0$. Taking a power is essential: nonzero $q(v)$ alone
does not assert nonzero $s(v)$ for that particular braid.

The braid $\gamma$ remains in $P_3^{(k)}$. Pairwise linking numbers
are homomorphisms on $P_3$, so they vanish on $P_3^{(1)}$ and every
positive derived subgroup. Its closure is algebraically split.
Lemma~\ref{lem:topology} proves $\gamma\notin\F_1^3$.
For $m>3$, add split vertical strands. This braid homomorphism preserves
derived membership, and its closure adds split unknots, preserving
ordinary signature. The same obstruction applies with $m$ meridians.
Finally $\F_n^m\subset\F_1^m$ for $n\geq1$
\cite[Section 2]{Otto}; set $k=n+2$.
\end{proof}

\section{An explicit example and verification scope}\label{sec:explicit}
With $a=\sigma_1^2$ and $b=\sigma_2^2$, put
\[
 c=[a,b],\quad d=a^{-1}ca,\quad e=b^{-1}cb,\quad
 u=[c,d],\quad v=[d,e],\quad \beta=[u,v].
\]
Normality of derived subgroups gives $c,d,e\in P_3^{(1)}$,
$u,v\in P_3^{(2)}$ and $\beta\in P_3^{(3)}$.
Its freely reduced Artin word has length $104$.
The accompanying package records a checkerboard calculation with a
$48\times48$ Goeritz matrix of inertia $(23,25,0)$ and zero correction,
and an independent Burau--Meyer calculation of signature $-2$.
The determinant of that Goeritz matrix is $-7154819319988224$.
These exact computations give a concrete $n=1$ witness; the all-depth
proof above does not depend on them.

The package preserves the original 44-file candidate, both independent
signature reconstructions, source locators and priority-search records.
The current release verifier authenticates its closed inventory and
replays specified finite checks; those checks supplement, rather than
prove, the universal topological and group-theoretic arguments.
Historical source-indexing and transport limitations are documented
alongside the corresponding corrections without rewriting the originals.

\paragraph{Attribution and historical scope.}
The signature quasimorphism and its values, amenable rigidity,
solvable-quotient defect restriction and signature obstruction are credited
inputs. The contribution here is their explicit connection to the
derived-pure-braid containment question, with an independently checked
finite witness. The source is the 2017 workshop report, published in
2018. Bounded primary-source searches through October 15, 2026 did not
establish an earlier complete answer to this precise question. This is
not a certificate of first priority or of continued openness.

\paragraph{AI assistance and review status.}
AI tools were used extensively in research, proof development, source
checking, computation, writing and adversarial verification. This note is
unrefereed. AI-assisted reviews do not constitute independent human or
journal peer review; no such review or formal proof-assistant
certification is claimed. The preserved initial candidate records one of five allowed proof
responses; the stronger argument is separately recorded as the second
response. Verification adds no further proof-search response.

\begin{thebibliography}{9}
\bibitem{Calegari} D.~Calegari, \emph{scl}, MSJ Memoirs \textbf{20},
Mathematical Society of Japan, 2009.
\href{https://doi.org/10.2969/msjmemoirs/020010000}{doi:10.2969/msjmemoirs/020010000}.

\bibitem{GG} J.-M.~Gambaudo and \'E.~Ghys, Braids and signatures,
\emph{Bull. Soc. Math. France} \textbf{133} (2005), 541--579.
\href{https://doi.org/10.24033/bsmf.2496}{doi:10.24033/bsmf.2496}.

\bibitem{HPR} S.~Harvey (joint work with J.~H.~Park and A.~Ray),
Pure braids, Whitney towers, and $0$-solvability, in
\emph{Mini-Workshop: Interactions between Low-dimensional Topology and Complex Algebraic Geometry}, Oberwolfach Reports \textbf{14} (2017),
report 50, pp.3020--3021, Question~3; published 2018.
\href{https://doi.org/10.4171/OWR/2017/50}{doi:10.4171/OWR/2017/50}.

\bibitem{KKMMM} M.~Kawasaki, M.~Kimura, S.~Maruyama,
T.~Matsushita and M.~Mimura, The space of nonextendable quasimorphisms,
\emph{Algebr. Geom. Topol.} \textbf{25} (2025), 1169--1226.
\href{https://doi.org/10.2140/agt.2025.25.1169}{doi:10.2140/agt.2025.25.1169}.

\bibitem{Otto} C.~Otto, The $(n)$-solvable filtration of link concordance
and Milnor's invariants, \emph{Algebr. Geom. Topol.} \textbf{14}
(2014), 2627--2654.
\href{https://doi.org/10.2140/agt.2014.14.2627}{doi:10.2140/agt.2014.14.2627}.

\bibitem{Toffoli} E.~Toffoli, The Atiyah--Patodi--Singer rho invariant
and signatures of links, \emph{Proc. Edinb. Math. Soc.} \textbf{65}
(2022), 404--440.
\href{https://doi.org/10.1017/S0013091522000153}{doi:10.1017/S0013091522000153}.
\end{thebibliography}
\end{document}
